\documentclass[aspectratio=169, 11pt]{beamer}

\usetheme{default}
\usecolortheme{dove}
\setbeamertemplate{navigation symbols}{}
\setbeamertemplate{footline}[frame number]
\setbeamercolor{frametitle}{fg=black}
\setbeamerfont{frametitle}{size=\large, series=\bfseries}

\usepackage[utf8]{inputenc}
\usepackage{graphicx}
\usepackage{booktabs}
\usepackage{xcolor}
\usepackage{hyperref}
\usepackage{array}
\usepackage{multirow}
\usepackage{tikz}
\usetikzlibrary{positioning, arrows.meta}
\graphicspath{{images/}}

\definecolor{linkblue}{RGB}{0,90,170}
\definecolor{dimgray}{RGB}{120,120,120}
\definecolor{goodgreen}{RGB}{34,139,34}
\definecolor{warnred}{RGB}{200,50,50}
\definecolor{noteorange}{RGB}{200,120,0}

\hypersetup{colorlinks=true, urlcolor=linkblue}

\title{ConSynth-X}
\subtitle{A Multi-Condition Synthetic Dataset for\\Construction Site Computer Vision Under Extreme Field Conditions}
\author{Viet Huy Duong \and Ruoxin Xiong, Ph.D.}
\institute{Kent State University}
\date{April 2026}

\begin{document}

% ═══════════════════════════════════════════════════════════════
% PART 1: PROBLEM & DATASET
% ═══════════════════════════════════════════════════════════════

\begin{frame}
\titlepage
\end{frame}

% ────────────────────────────────────────
\begin{frame}{The Problem}

\textbf{Goal:} Provide a comprehensive dataset for construction site computer vision that enables \textbf{training and evaluating} models across diverse extreme field conditions.

\vspace{0.6em}
\textbf{Why this matters:}
\begin{itemize}\setlength\itemsep{0.15em}
    \item Construction sites face rain, snow, fog, low light, glare
    \item Models trained on clean data fail dramatically in the field:
\end{itemize}

\vspace{0.2em}
\centering
\footnotesize
\begin{tabular}{lcc}
\toprule
\textbf{Condition} & \textbf{Detection mAP@0.5} & \textbf{Drop} \\
\midrule
Clear weather (original) & 0.819 & --- \\
Rain / Snow & 0.446 & {\color{warnred}$-$46\%} \\
Nighttime & 0.454 & {\color{warnred}$-$45\%} \\
\bottomrule
\end{tabular}

\vspace{0.6em}
\raggedright
\begin{itemize}\setlength\itemsep{0.15em}
    \item No existing dataset covers multiple extreme conditions for construction
    \item Collecting real extreme-condition data is dangerous and expensive
\end{itemize}

\vspace{0.3em}
$\to$ We synthetically generate a multi-condition dataset to fill this gap.
\end{frame}

% ────────────────────────────────────────
\begin{frame}{Our Solution: ConSynth-X}

\textbf{We take existing construction photos and digitally add weather, night, fog, etc.}

\vspace{0.5em}
\centering
\footnotesize
\begin{tabular}{llcl}
\toprule
\textbf{What we add} & \textbf{How} & \textbf{Images} & \textbf{In simple terms} \\
\midrule
Rain / Snow (style) & Neural Style Transfer & 57K & Copy weather ``look'' from reference photos \\
Rain / Snow (diff.) & InstructPix2Pix AI & 10K & AI edits image based on text prompt \\
Nighttime & CycleGAN-Turbo & 23K & AI converts day photo to night \\
Fog & Koschmieder model & 3K & Physics-based fog simulation \\
Smaller objects & FLUX.1 outpainting & 2K & Expand image canvas, objects appear farther \\
\midrule
\multicolumn{2}{l}{\textbf{Total augmented}} & \textbf{95K+} & \\
\bottomrule
\end{tabular}

\vspace{0.5em}
\raggedright
{\small Applied to 2 datasets: \textbf{Construction Site 10k} (3 classes) + \textbf{SODA} (15 classes).\\
All original bounding box annotations are preserved in augmented images.}
\end{frame}

% ────────────────────────────────────────
\begin{frame}{Visual Examples}
\centering
\includegraphics[width=\textwidth, height=0.78\textheight, keepaspectratio]{fig1_comparison_grid.png}

\vspace{0.2em}
{\small Each row: same photo $\to$ Style Transfer $\to$ AI Diffusion + physics rain/snow overlay.}
\end{frame}

% ────────────────────────────────────────
\begin{frame}{Prior Work \& Our Contribution}

\begin{columns}[T]
\begin{column}{0.48\textwidth}
\textbf{Ding et al.\ (2024):}
\begin{itemize}\setlength\itemsep{0.1em}
    \item UIA-YOLOv5 + Neural Style Transfer
    \item ExtCon: 506 real extreme images
    \item +8.21\% mAP improvement
    \item Only 1 method (NST), no snow/night/scale
\end{itemize}
\end{column}
\begin{column}{0.48\textwidth}
\textbf{ConSynth-X (ours):}
\begin{itemize}\setlength\itemsep{0.1em}
    \item 126,000+ images (250$\times$ larger)
    \item 3 complementary pipelines
    \item 5 condition categories
    \item 5-dimensional validation framework
    \item Open-source pipeline for new datasets
\end{itemize}
\end{column}
\end{columns}

\vspace{0.8em}
\textbf{Key question:} How do we know the synthetic images are good enough?\\
$\to$ We validate from 5 different angles.
\end{frame}

% ═══════════════════════════════════════════════════════════════
% PART 2: VALIDATION (5 approaches)
% ═══════════════════════════════════════════════════════════════

% ────────────────────────────────────────
\begin{frame}{How We Validate: 5 Questions}

{\small Each approach answers a different question about image quality:}

\vspace{0.5em}
\begin{enumerate}\setlength\itemsep{0.3em}
    \item \textbf{Weather Recognizability} --- Can you tell it's raining/snowing?
    \item \textbf{Distributional Fidelity} --- Is it statistically similar to real weather photos?
    \item \textbf{Texture Fidelity} --- Are surface details (concrete, metal, fabric) preserved?
    \item \textbf{Embedding Distance} --- In AI ``feature space,'' is it close to real weather?
\end{enumerate}

\vspace{0.3em}
{\small \textbf{+ Supplementary:} AI-generated image detection (UnivFD) --- confirms our images retain real-photo statistics, but this metric is \textbf{not informative} for edit-based augmentation (see backup slide).}

\vspace{0.3em}
No single metric tells the whole story.\\
\textbf{Cross-validation} across all 4 reveals trade-offs between pipelines.
\end{frame}

% ────────────────────────────────────────
\begin{frame}{Validation 1: Weather Recognizability --- Can You Tell It's Raining?}

\begin{columns}[T]
\begin{column}{0.52\textwidth}
\textbf{Question:} Is the weather effect actually visible?\\
{\small (An image can look ``real'' but not look like ``rain.'')}

\vspace{0.2em}
\textbf{Method:} Pre-trained \textbf{weather classifier} (SigLIP2) classifies each augmented image into 5 weather classes.

\vspace{0.2em}
\textbf{Caveat:} Classifier trained on driving scenes --- it misclassifies 58\% of \textit{original clear} construction images as ``rain'' (cross-domain bias). \textbf{Relative comparison} between pipelines is more meaningful than absolute numbers.
\end{column}
\begin{column}{0.46\textwidth}
\centering
\footnotesize
\begin{tabular}{lc}
\toprule
\textbf{Condition} & \textbf{Accuracy} \\
\midrule
Original (clear) & 6\% sun \\
\midrule
ST Rain A & \textbf{89\%} \\
ST Snow B & \textbf{95\%} \\
\midrule
Diff.\ Rain & 64\% \\
Diff.\ Snow & {\color{warnred}26\%} \\
\bottomrule
\end{tabular}

\vspace{0.5em}
\textbf{Takeaway:}\\
Style transfer $\to$ strong, obvious weather.\\
Diffusion $\to$ subtle, less recognizable.\\[0.3em]
{\color{warnred}Diffusion snow only 26\%}, but its texture quality is the best (see Validation 4).
\end{column}
\end{columns}
\end{frame}

% ────────────────────────────────────────
\begin{frame}{Validation 2: FID/KID --- Is It Statistically Similar to Real Weather?}

\begin{columns}[T]
\begin{column}{0.55\textwidth}
\textbf{Question:} How close is the distribution of our synthetic images to the distribution of \textit{real} weather photos?

\vspace{0.3em}
\textbf{Method:} \textbf{FID} (Fr\'{e}chet Inception Distance, Heusel et al.\ 2017) compares feature distributions between two image sets. Lower = more similar.

\vspace{0.3em}
\textbf{How:} Compare augmented images against 3 real-weather datasets (ACDC, WeatherBench, WeatherNet).

\vspace{0.2em}
``Baseline gap'' = how far original clear images are from real weather. If augmentation \textit{reduces} this gap, it's working.
\end{column}
\begin{column}{0.43\textwidth}
\centering
\footnotesize
\begin{tabular}{lcc}
\toprule
 & \textbf{FID} & $\boldsymbol{\Delta}$ \\
\midrule
\multicolumn{3}{l}{\textit{Night (vs ACDC night):}} \\
Original & 270.5 & --- \\
CycleGAN & \textbf{178.9} & {\color{goodgreen}$-$34\%} \\
\midrule
\multicolumn{3}{l}{\textit{Snow (vs ACDC snow):}} \\
Original & 216.0 & --- \\
Diffusion & \textbf{205.3} & {\color{goodgreen}$-$5\%} \\
\midrule
\multicolumn{3}{l}{\textit{Rain (vs ACDC rain):}} \\
Original & 189.8 & --- \\
ST rain & 192.3 & $\approx$0\% \\
\bottomrule
\end{tabular}

\vspace{0.3em}
\textbf{Night is the biggest win:}\\
34\% closer to real night photos.
\end{column}
\end{columns}
\end{frame}

% ────────────────────────────────────────
\begin{frame}{Validation 3: Texture Fidelity --- Are Surface Details Preserved?}

\begin{columns}[T]
\begin{column}{0.55\textwidth}
\textbf{Question:} Does augmentation damage the fine texture of surfaces (concrete, metal, fabric)?

\vspace{0.3em}
\textbf{Method:} Adapted from \textbf{Duminil et al.\ (2025, Sci.\ Reports)}. Extract 4 texture features per image (GLCM, LBP, DCT, Haralick), compare distributions via Wasserstein distance, then fuse with \textbf{Dempster-Shafer belief theory}: all 4 features ``vote'' on texture quality.

\vspace{0.2em}
{\footnotesize $H$ = belief ``faithful'' (preserved). $\bar{H}$ = belief ``artifact'' (damaged).}
\end{column}
\begin{column}{0.43\textwidth}
\centering
\footnotesize
\begin{tabular}{lccl}
\toprule
 & $H$ & $\bar{H}$ & \\
\midrule
{\color{goodgreen}Diff.\ Snow} & {\color{goodgreen}\textbf{.86}} & .00 & Faithful \\
{\color{goodgreen}Diff.\ Rain} & {\color{goodgreen}\textbf{.85}} & .00 & Faithful \\
\midrule
ST Rain & .00 & .26 & Uncertain \\
ST Snow & .00 & .83 & Artifact \\
\midrule
Fog & .00 & .83 & Artifact \\
Night & .02 & .73 & Artifact \\
\bottomrule
\end{tabular}

\vspace{0.5em}
\textbf{Diffusion preserves texture}\\(all 4 criteria agree: $H{>}0.85$).

\vspace{0.3em}
Style transfer creates visible\\weather but damages texture.

\vspace{0.3em}
Night has highest ``conflict'':\\local texture OK, global lost.
\end{column}
\end{columns}
\end{frame}

% ────────────────────────────────────────
\begin{frame}{Validation 4: Embedding Distance --- The Idea}

\textbf{Question:} In AI ``feature space,'' are our augmented images closer to \textit{real} weather photos?

\vspace{0.3em}
\textbf{Method:} Relative Mahalanobis Distance (Ruck et al.\ 2026). Subtract background to isolate weather signal.

\vspace{0.2em}
\centering
\begin{tikzpicture}[
    every node/.style={font=\footnotesize},
    blob/.style={ellipse, draw, minimum width=2.2cm, minimum height=1cm, align=center},
    >={Stealth[length=2mm]}
]
% Number line
\draw[->, thick, gray!50] (-0.5,0) -- (11.5,0);
\node[below, font=\scriptsize, gray] at (11.5,0) {distance to real weather};

% Real weather (ACDC) cluster
\node[blob, fill=goodgreen!15, draw=goodgreen] (real) at (1.5,0.8) {Real weather\\(ACDC)};
\draw[goodgreen, thick] (1.5,0) -- (1.5,0.3);
\node[below, font=\scriptsize\bfseries, goodgreen] at (1.5,-0.2) {$-d_{rel} \approx 0$};

% Augmented
\node[blob, fill=linkblue!15, draw=linkblue] (aug) at (5.5,0.8) {Augmented\\(ours)};
\draw[linkblue, thick] (5.5,0) -- (5.5,0.3);
\node[below, font=\scriptsize\bfseries, linkblue] at (5.5,-0.2) {$-d_{rel} = -25.6$};

% Original clear
\node[blob, fill=warnred!15, draw=warnred] (orig) at (9.5,0.8) {Original\\(clear day)};
\draw[warnred, thick] (9.5,0) -- (9.5,0.3);
\node[below, font=\scriptsize\bfseries, warnred] at (9.5,-0.2) {$-d_{rel} = -39.3$};

% Gap arrows
\draw[<->, thick, goodgreen!70] (1.7,-0.7) -- (5.3,-0.7);
\node[below, font=\scriptsize, goodgreen!80] at (3.5,-0.7) {remaining gap};

\draw[<->, thick, linkblue!70] (5.7,-0.7) -- (9.3,-0.7);
\node[below, font=\scriptsize\bfseries, linkblue] at (7.5,-0.7) {gap closed = \textbf{38\%}};

\end{tikzpicture}

\vspace{0.2em}
\raggedright
{\small \textbf{Example:} Night (CLIP). Original is $-39.3$ from real ACDC night. CycleGAN augmentation moves it to $-25.6$ --- closing 38\% of the gap.}

\vspace{0.2em}
{\small \textbf{``\% gap closed''} $= \dfrac{\text{original} - \text{augmented}}{\text{original} - \text{baseline}} \times 100$}

\end{frame}

% ────────────────────────────────────────
\begin{frame}{Validation 4: Results \& Cross-Domain Limitation}

\begin{columns}[T]
\begin{column}{0.48\textwidth}
\centering
\footnotesize
\textbf{CLIP (semantic/style)}

\vspace{0.3em}
\begin{tabular}{lcc}
\toprule
\textbf{Condition} & \textbf{Best method} & \textbf{Gap closed} \\
\midrule
Night & CycleGAN & \textbf{38\%} \\
Snow & ST heavy & 27\% \\
Rain & ST moderate & 25\% \\
Fog & Physics & 18\% \\
\bottomrule
\end{tabular}

\vspace{0.5em}
{\color{goodgreen}\textbf{All conditions move closer}}\\
{\color{goodgreen}to real weather in CLIP space.}

\vspace{0.3em}
{\small CLIP captures ``weather look'' (style/semantics) which is relatively \textbf{domain-agnostic}: rain looks like rain whether on a road or a construction site.}
\end{column}
\begin{column}{0.48\textwidth}
\centering
\footnotesize
\textbf{DINOv3 (texture/structure)}

\vspace{0.3em}
\begin{tabular}{lcc}
\toprule
\textbf{Condition} & \textbf{Best method} & \textbf{Gap closed} \\
\midrule
Night & CycleGAN & 9\% \\
Snow & IP2P & 5\% \\
Rain & ST moderate & 1\% \\
Fog & Physics & 1\% \\
\bottomrule
\end{tabular}

\vspace{0.5em}
{\color{warnred}\textbf{Near-zero movement.}}\\
{\color{warnred}But this is a \textbf{metric limitation}:}

\vspace{0.3em}
{\small DINOv3 captures texture/structure. Reference is ACDC (\textit{driving} scenes) --- wet roads $\neq$ wet concrete. The \textbf{cross-domain structural gap} overwhelms the weather signal.}
\end{column}
\end{columns}

\vspace{0.5em}
{\small \textbf{Takeaway:} CLIP results validate our augmentation. DINOv3 gap reflects \textbf{domain mismatch} (driving vs.\ construction), not augmentation failure. A construction-specific weather reference would be needed --- but that is exactly what does not exist yet.}
\end{frame}

% ═══════════════════════════════════════════════════════════════
% PART 3: SYNTHESIS & CONCLUSIONS
% ═══════════════════════════════════════════════════════════════

% ────────────────────────────────────────
\begin{frame}{Putting It All Together: The Trade-Off}

\centering
\footnotesize
\begin{tabular}{lcccc}
\toprule
\textbf{What we measured} & \textbf{Style Transfer} & \textbf{Diffusion} & \textbf{Night} \\
\midrule
1. Weather visible? (classifier $\uparrow$) & \textbf{62--95\%} & 26--64\% & N/A \\
2. Close to real weather? (FID $\downarrow$) & 192--228 & 194--219 & \textbf{179} \\
3. Texture preserved? ($H$ $\uparrow$) & 0.00 & \textbf{0.86} & 0.02 \\
4. Embedding shift? (CLIP $\uparrow$) & 25--27\% & 18\% & \textbf{38\%} \\
\bottomrule
\end{tabular}

\vspace{0.8em}
\raggedright
\begin{columns}[T]
\begin{column}{0.48\textwidth}
\textbf{The key finding:}\\[0.2em]
Style Transfer and Diffusion are \textbf{complementary}:
\begin{itemize}\setlength\itemsep{0.1em}
    \item Style Transfer: strong weather, damages texture
    \item Diffusion: natural texture, subtle weather
\end{itemize}
Neither is ``better'' --- they serve different purposes.
\end{column}
\begin{column}{0.48\textwidth}
\textbf{Practical recommendation:}
\begin{itemize}\setlength\itemsep{0.1em}
    \item \textbf{Training data:} use Diffusion\\(natural texture helps model learn)
    \item \textbf{Robustness testing:} use Style Transfer\\(extreme weather stress-tests the model)
    \item \textbf{Night:} best overall domain shift
\end{itemize}
\end{column}
\end{columns}
\end{frame}

% ────────────────────────────────────────
\begin{frame}{Downstream Detection: Does Augmented Data Help?}

\begin{columns}[T]
\begin{column}{0.52\textwidth}
\textbf{Question:} If we add synthetic weather images to training, does it improve or hurt the model?

\vspace{0.3em}
\textbf{Setup:} YOLOv8n detector trained with/without augmented data, tested on clean images.

\vspace{0.3em}
\centering
\footnotesize
\begin{tabular}{lcc}
\toprule
\textbf{Training data} & \textbf{mAP@.5} & \textbf{mAP@.5:.95} \\
\midrule
Original only (7K) & 0.562 & 0.338 \\
+Weather (17K) & 0.547 & 0.339 \\
+Weather strict & 0.538 & \textbf{0.349} \\
\bottomrule
\end{tabular}
\end{column}
\begin{column}{0.46\textwidth}
\vspace{1em}
\textbf{Result:} Adding 10K+ synthetic images does \textbf{not degrade} clean-condition detection (mAP@.5:.95 neutral $\pm$0.001).

\vspace{0.5em}
This means synthetic quality is sufficient for training --- the model doesn't get ``confused'' by augmented images.

\vspace{0.5em}
{\color{noteorange}\textbf{Next step:}} Test on weather/night conditions (experiment ready, awaiting compute).\\
This will show the \textit{robustness benefit}.
\end{column}
\end{columns}
\end{frame}

% ────────────────────────────────────────
\begin{frame}{Paper Status \& Next Steps}

\begin{columns}[T]
\begin{column}{0.48\textwidth}
\textbf{Target:} Nature Scientific Data

\vspace{0.5em}
\textbf{Completed:}
\begin{itemize}\setlength\itemsep{0.1em}
    \item 5 augmentation pipelines
    \item 5-dimensional validation
    \item 23-page manuscript, 21 references
    \item Prior work comparison (Ding et al.)
    \item Cross-validation synthesis
\end{itemize}

\vspace{0.5em}
\textbf{Review feedback:}\\
{\small ``Underlying research has merit... could reach Minor Revision quality''}
\end{column}
\begin{column}{0.48\textwidth}
\textbf{Remaining:}
\begin{enumerate}\setlength\itemsep{0.1em}
    \item Cross-condition detection results\\{\footnotesize (experiment script ready, $\sim$2h GPU)}
    \item SODA 15-class validation
    \item Dataset upload (HuggingFace)
    \item Final paper polish
\end{enumerate}

\vspace{0.5em}
\textbf{Roadmap:}\\
{\small Paper 1 (Dataset) $\to$ Paper 2 (VLM Benchmark)\\$\to$ Paper 3 (Detection Improvement)}
\end{column}
\end{columns}
\end{frame}

% ═══════════════════════════════════════════════════════════════
% BACKUP SLIDES
% ═══════════════════════════════════════════════════════════════

\appendix

% ────────────────────────────────────────
\begin{frame}{Backup: Why Not Use AI-Generated Image Detectors?}

\begin{columns}[T]
\begin{column}{0.48\textwidth}
\textbf{Two different detection problems:}

\vspace{0.3em}
\textbf{SID} (Synthetic Image Detection):\\
{\small Classify \textit{entire} image as real vs AI-generated.\\
Examples: UnivFD, DIRE, NPR.\\
Designed for: Midjourney, DALL-E, StyleGAN outputs.}

\vspace{0.3em}
\textbf{IML} (Image Manipulation Localization):\\
{\small Detect \textit{which regions} of a real image were edited.\\
Examples: TruFor, X-Edit, DiffForensics.\\
Designed for: splicing, inpainting, local edits.}
\end{column}
\begin{column}{0.48\textwidth}
\textbf{Our images are neither:}

\vspace{0.3em}
ConSynth-X starts from \textbf{real photos} and applies \textbf{global edits} (weather, night).

\vspace{0.3em}
$\to$ SID methods see ``real photo'' (99.5\% fooling)\\because most pixels are original.

\vspace{0.3em}
$\to$ IML methods expect \textbf{local} edits,\\but weather is a \textbf{global} transformation.

\vspace{0.5em}
\textbf{Our approach:} Instead of binary real/fake, we validate through \textbf{domain-specific metrics} --- weather recognition, distributional fidelity, texture preservation, and embedding distance.
\end{column}
\end{columns}

\vspace{0.3em}
{\footnotesize UnivFD results ($>$99.5\% fooling) are reported as supplementary. Full results in paper appendix.}
\end{frame}

% ────────────────────────────────────────
\begin{frame}{}
\centering
\vspace{2cm}
{\Huge\bfseries Thank You}

\vspace{1cm}
{\large Questions?}

\vspace{1cm}
{\footnotesize
Viet Huy Duong --- \texttt{vduong1@kent.edu}\\
Ruoxin Xiong --- \texttt{rxiong3@kent.edu}\\
Kent State University
}
\end{frame}

\end{document}
